\section*{Chapter \ref{ch:b3:total-synthesis} --- Total Synthesis: Strategy and Classic Routes}

\begin{solution}{exo:b3:total-synthesis:1}
At 90\,\%: $0.9^{10} = 35$\,\% and $0.9^{20} = 12$\,\%. At 80\,\%: $0.8^{10} = 11$\,\% and $0.8^{20} = 1.2$\,\%.
\end{solution}

\begin{solution}{exo:b3:total-synthesis:2}
$5 + 3 + 1 + 4 = 13$ steps in all; LLS $5 + 1 + 4 = 10$.
\end{solution}

\begin{solution}{exo:b3:total-synthesis:3}
4-(Dimethylamino)butan-2-one, $\ce{CH3COCH2CH2N(CH3)2}$: the enol of acetone adds to the iminium
ion $\ce{CH2=N+(CH3)2}$.
\end{solution}

\begin{solution}{exo:b3:total-synthesis:4}
Disconnect the two ring bonds allylic to the C=C: buta-1,3-diene and but-3-en-2-one
(methyl vinyl ketone), which react on heating.
\end{solution}

\begin{solution}{exo:b3:total-synthesis:5}
Linear: $0.9^{21} = 11$\,\%. Convergent: $0.9^{11} = 31$\,\%, almost three times as much.
\end{solution}

\begin{solution}{exo:b3:total-synthesis:6}
A: $(7 + 1)/12 = 67$\,\%. B: $7/9 = 78$\,\%.
\end{solution}

\begin{solution}{exo:b3:total-synthesis:7}
Without: $0.9^8 = 43$\,\%. With four extra steps at 95\,\%: $43\,\% \times 0.95^4 = 35$\,\%.
\end{solution}

\begin{solution}{exo:b3:total-synthesis:8}
Primary: a bulky silyl ether (\textit{tert}-butyldiphenylsilyl), removed with fluoride.
Secondary: a 4-methoxybenzyl ether, removed by oxidation with DDQ. The tertiary
alcohol, hindered, can often be left free, or protected as a small silyl ether
cleaved by mild acid. Each condition leaves the other groups untouched.
\end{solution}

\begin{solution}{exo:b3:total-synthesis:9}
The total syntheses needed tens of steps with \omterm{def:b3:total-synthesis:architecture}{overall yields} far below 1\,\%; a
compound with the full ring system and most stereocentres, extracted from the
needles of the European yew, a renewable source, is converted in a few steps.
\end{solution}

\begin{solution}{exo:b3:total-synthesis:10}
Michael addition: the enolate of the dione (C2, between the carbonyls) adds to the
$\beta$-carbon of methyl vinyl ketone, giving a triketone. Aldol: the methyl ketone's
enolate attacks a ring carbonyl intramolecularly, and dehydration gives the
cyclohexenone: the Wieland--Miescher ketone, a bicyclic enedione with an angular
methyl group.
\end{solution}

\begin{solution}{exo:b3:total-synthesis:11}
$y^{m+1}/y^{2m+1} = y^{-m}$: at 90\,\%, 1.7 for $m = 5$ and 4.9 for $m = 15$. The longer the route, the more
convergence pays.
\end{solution}

\begin{solution}{exo:b3:total-synthesis:12}
Each ring closes with the cation and the attacking double bond in a chair-like
arrangement; addition across each double bond is anti, and the substituents take
pseudo-equatorial positions, which makes every ring fusion \textit{trans}. An \textit{E} double bond
places its two chain continuations so that the \textit{trans} fusion results; a \textit{Z} one would
give the \textit{cis} fusion: the geometry of the polyene encodes the stereochemistry of
the product.
\end{solution}

\begin{solution}{pb:b3:total-synthesis:1}
\textbf{1.} $\ce{C4H6O2}$, $\ce{CH5N}$, $\ce{C5H6O5}$, $\ce{C8H13NO}$.

\textbf{2.} $\ce{C4H6O2 + CH5N + C5H6O5 -> C8H13NO + 2CO2 + 2H2O}$.

\textbf{3.} 86.0, 31.0 and \qty{146.0}{g/mol}; tropinone \qty{139.0}{g/mol}.

\textbf{4.} $139.0/(86.0 + 31.0 + 146.0) = 0.53$: 53\,\%.

\textbf{5.} Its central $\ce{CH2}$ groups, each between two carbonyl groups, enolise readily in
water; acetone itself is too weakly nucleophilic, and the activating carboxyl
groups leave afterwards.

\textbf{6.} The iminium ions need free amine (too much acid protonates it), and the
enol and the iminium formation need some acid: a pH near neutrality balances them.

\textbf{7.} $\ce{CH3NH2}$ adds to one aldehyde group; loss of water gives the iminium ion
$\ce{R-CH=N+H-CH3}$.

\textbf{8.} An enol carbon of acetonedicarboxylic acid attacks the iminium carbon, making a
C--C bond and a secondary amine.

\textbf{9.} The secondary amine condenses with the second aldehyde group of the same
chain, giving a cyclic (five-membered) iminium ion.

\textbf{10.} The enol on the other side of the ketone attacks it, closing the
six-membered ring and the bicyclic skeleton.

\textbf{11.} Each is a $\beta$-keto acid, which loses $\ce{CO2}$ through a cyclic six-membered transition
state, giving the enol of the ketone.

\textbf{12.} Two C--C bonds and two C--N bonds.

\textbf{13.} $0.75^{14} = 0.018$: 1.8\,\%.

\textbf{14.} $42/1.78 = 24$.

\textbf{15.} 100\,\%: one step, all bond construction.

\textbf{16.} One operation, water as solvent, cheap components, no protecting groups and
almost no waste besides carbon dioxide and water.

\textbf{17.} The plant builds the tropane skeleton from an amino acid-derived
cyclic iminium ion and an acetate-derived unit, through the same kind of Mannich
chemistry.

\textbf{18.} Cascades (and ring-forming \omterm{def:b3:pericyclic:families}{cycloadditions} such as the Diels--Alder
reaction).

\textbf{19.} No: a \omterm{def:b3:point-groups:mirror}{mirror plane} passes through N, its methyl group, C3 and the carbonyl
oxygen.

\textbf{20.} Each has four different groups, but they are mirror images of each other in
the molecule's own plane of symmetry: a meso-type compound.

\textbf{21.} Diastereomers, differing in the orientation of the OH on C3 relative to
the nitrogen bridge.

\textbf{22.} No: the \omterm{def:b3:point-groups:mirror}{mirror plane} is kept in both.

\textbf{23.} The two faces of the carbonyl group are not equivalent: one faces the
nitrogen bridge, the other the two-carbon bridge, which hinder differently.

\textbf{24.} The one-pot route gives about 24 times the \omterm{def:b3:total-synthesis:architecture}{overall yield} of the 14-step
linear route.
\end{solution}
